\documentclass[11pt]{article}
\usepackage{preamble}
\setlength{\parskip}{4pt}

\begin{document}

\daybanner{AP Statistics \textperiodcentered\ Unit 4 \textperiodcentered\ Topic 4.5 \textperiodcentered\ B: 2027-01-21 \textperiodcentered\ E: 2027-03-08 \textperiodcentered\ TEACHER KEY}{A Test of Mean Resolution Time}

\sectionbanner{CED TETHER}

{\small
\begin{itemize}[leftmargin=1.5em,itemsep=2pt,topsep=2pt]
  \item \textbf{Skill 3.E} --- Calculate a test statistic and find a p-value, provided conditions for inference are met.
  \item \textbf{Skill 4.B} --- Interpret statistical calculations and findings to assign meaning or assess a claim.
  \item \textbf{Skill 4.E} --- Justify a claim using a decision based on significance tests.
  \item \textbf{EU VAR-7} --- The t-distribution may be used to model variation.
  \item \textbf{LO VAR-7.E} --- Calculate an appropriate test statistic for a population mean, including the mean difference between values in matched pairs. [Skill 3.E]
  \item \textbf{EK VAR-7.E.1} --- For a single quantitative variable when random sampling with replacement from a population that can be modeled with a normal distribution with mean $\mu$ and standard deviation $\sigma$, the sampling distribution of t = ($\bar{x}$ - $\mu$) / (s/$\surd$n) has a t-distribution with n - 1 degrees of freedom.
  \item \textbf{EU DAT-3} --- Significance testing allows us to make decisions about hypotheses within a particular context.
  \item \textbf{LO DAT-3.E} --- Interpret the p-value of a significance test for a population mean, including the mean difference between values in matched pairs. [Skill 4.B]
  \item \textbf{EK DAT-3.E.1} --- An interpretation of the p-value of a significance test for a population mean should recognize that the p-value is computed by assuming that the null hypothesis is true, i.e., by assuming that the true population mean is equal to the particular value stated in the null hypothesis.
  \item \textbf{LO DAT-3.F} --- Justify a claim about the population based on the results of a significance test for a population mean. [Skill 4.E]
  \item \textbf{EK DAT-3.F.1} --- A formal decision explicitly compares the p-value to the significance $\alpha$. If the p-value $\leq$ $\alpha$, then reject the null hypothesis, $H_0$: $\mu$ = $\mu_0$. If the p-value > $\alpha$, then fail to reject the null hypothesis.
  \item \textbf{EK DAT-3.F.2} --- The results of a significance test for a population mean can serve as the statistical reasoning to support the answer to a research question about the population that was sampled.
\end{itemize}
}
\textbf{Essential question:} What conclusion is warranted when a one-sample t-test does not cross its evidence threshold?\par
\textbf{DOK-3 skill:} 4.E \textperiodcentered\ \textbf{FRQ pattern:} {\small\texttt{mean-\allowbreak{}test-\allowbreak{}decision-\allowbreak{}and-\allowbreak{}conclusion}}\par
\textbf{DOK rationale:} Part (c) coordinates the lower-tail p-value, preset threshold, asymmetric decision logic, and SRS population scope to adjudicate two conclusions and trace the consequence of overclaiming.\par

\frameworkphaseheader{Do Now (first take) $\rightarrow$ Explore (video + follow-along) $\rightarrow$ Exit (finish + turn in)}{1 $\rightarrow$ 3}{45}{%
\begin{itemize}[leftmargin=1.5em,itemsep=2pt,topsep=2pt]
  \item Display the board and ask for one sentence using the study description.
  \item Begin the 7.5 video follow-along at minute 5.
  \item Return to the board at minute 33. Students finish the shortened parts and justify their judgment.
  \item Collect at the door. Score part (c) only using the three required elements.
\end{itemize}
}{%
\begin{itemize}[leftmargin=1.5em,itemsep=2pt,topsep=2pt]
  \item Write a one-sentence first take before the video.
  \item Complete the video follow-along worksheet.
  \item Finish (a)--(c), revise the first take in the exit line, and turn in the sheet.
\end{itemize}
}{%
\begin{itemize}[leftmargin=1.5em,itemsep=2pt,topsep=2pt]
  \item Which specific number or study detail supports your choice?
  \item What claim would go beyond the evidence, and why?
\end{itemize}
}{Circulate and ask for evidence. Accept a concise response; do not require omitted calculations or a second full inference write-up.}

\begin{callout}{\IconWarn\ WATCH FOR}{calloutred}
\begin{itemize}[leftmargin=1.5em,itemsep=2pt,topsep=2pt]
  \item Choosing a speaker without a reason.
  \item Copying numbers without explaining what they support.
\end{itemize}
\end{callout}

\sectionbanner{SCORING PART (c) --- E / P / I}

\textbf{Expected elements}
\begin{itemize}[leftmargin=1.5em,itemsep=2pt,topsep=2pt]
  \item \textbf{decision} (required) --- Uses 0.0645 > 0.05 to fail to reject
  \item \textbf{context} (required) --- States insufficient evidence for a mean below 12 hours among the 1,200 tickets
  \item \textbf{limits} (required) --- Explains why the test does not prove equality or chance as the only explanation
\end{itemize}

{\small\renewcommand{\arraystretch}{1.25}
\noindent\begin{tabular}{|p{0.5in}|p{5.9in}|}
\hline
\textbf{E} & Uses the threshold correctly, gives the population conclusion, and explains both limits on what the test proves. \\ \hline
\textbf{P} & Shows substantive valid reasoning for at least one required element, but does not meet all E requirements. Credit correct reasoning even when another claim is wrong. \\ \hline
\textbf{I} & Shows no substantive valid reasoning for any required element; a name, unsupported choice, or copied number alone does not count. \\ \hline
\end{tabular}}

\clearpage
\focusbanner{TODAY'S PROBLEM --- annotated}

Suppose a help desk closed 1,200 tickets in a quarter. To investigate whether their population mean resolution time was below 12 hours, an analyst took an SRS of 24 and used $\alpha=0.05$.\par\smallskip \textbf{Leah:} ``The lower-tail test gives $t=-1.5747$, $df=23$, and $p\approx0.0645$. I would state what evidence this does and does not provide.''\par \textbf{Damon:} ``Because $p>0.05$, accept $H_0$. The test proves the population mean was exactly 12 hours and the observed difference was only chance.''\par
\begin{center}
\textbf{Ticket sample}\par\smallskip
\begin{tabular}{|l|c|}
\hline
\textbf{Summary} & \textbf{Value} \\ \hline
\textbf{Population} & 1200 \\ \hline
\textbf{$n$} & 24 \\ \hline
\textbf{Mean (h)} & 11.1 \\ \hline
\textbf{SD (h)} & 2.8 \\ \hline
\textbf{Shape} & Symmetric; no outliers \\ \hline
\end{tabular}
\end{center}
\begin{firsttakebox}{FIRST TAKE (5 min)}
Does this small study seem able to prove that the average was exactly 12 hours? Explain in one sentence.\par
\answer{Ungraded. Everyday language is enough before the video; formal vocabulary belongs in the finish.}
\end{firsttakebox}

\textit{Rules callout ``COMPARE EVIDENCE TO THE PRESET THRESHOLD (keep on the page)'' is printed on the student sheet.}\par\medskip

\dokbadge{1}~\textbf{(a)} State $H_0$ and $H_a$ for the population mean resolution time.\par
\answer{$H_0:\mu=12$ hours; $H_a:\mu<12$ hours, where $\mu$ is the mean for all 1,200 tickets closed that quarter.}

\dokbadge{2}~\textbf{(b)} Using $\bar x=11.1$, $s=2.8$, and $n=24$, calculate $t$. Use the supplied $p=0.0645$; you do not need to find it again.\par
\answer{$t=(11.1-12)/(2.8/\sqrt{24})=-1.5747$. The supplied lower-tail $p$-value is about 0.0645 with $df=23$.}

\dokbadge{3}~\textbf{(c)} In 3--4 sentences, decide whose account is justified. Compare the p-value with 0.05, state the conclusion for all 1,200 tickets, and explain why failure to reject does not prove equality or that chance was the only explanation.\par
\answer{Leah is justified: $0.0645>0.05$, so fail to reject $H_0$. There is not convincing evidence that the mean resolution time for all 1,200 tickets was below 12 hours. The p-value measures how unusual the data would be assuming a mean of 12; it does not prove that assumption. Chance remains a possible explanation, but the test does not prove it was the only explanation.}

\begin{summaryexitbox}{EXIT}
One thing the video changed about my first take: \hrulefill
\end{summaryexitbox}

\end{document}
